\chapter{Аналитическая часть}\label{ch:-4}

В данном разделе будут рассмотрены теоретические аспекты алгоритмов нахождения заданного значения в словаре.


\section{Линейный поиск}

Линейный поиск проходит через массив, начиная с первого элемента, и сравнивает каждый элемент с искомым.
Если элемент найден, алгоритм возвращает его индекс.
Если нет — продолжается поиск до конца массива.
В худшем случае требуется проверить все элементы массива.
Элементы, расположенные в начале массива будут найдены быстрее тех, что расположены в конце.


Пусть массив $A = \{a_1, a_2, \dots, a_n\}$, а искомый элемент $x$. Тогда алгоритм линейного поиска можно записать следующим образом:

\begin{itemize}
    \item для каждого элемента $a_i \in A$, если $a_i = x$, то вернуть индекс $i$;
    \item если $x$ не найден после проверки всех элементов, вернуть $-1$.
\end{itemize}


\section{Бинарный поиск}

Бинарный поиск работает на отсортированных массивах.
Алгоритм разделяет массив пополам и проверяет центральный элемент.
Если центральный элемент совпадает с искомым, то поиск завершён.
Если искомый элемент меньше центрального, алгоритм продолжает работу в левой половине массива, если больше — в правой.
Этот процесс повторяется до тех пор, пока не будет найден элемент или интервал для поиска не станет пустым.

\section*{Вывод}

В данном разделе были рассмотрены теоретические аспекты алгоритмов нахождения заданного значения в словаре.

